\DocumentMetadata{
  pdfversion=2.0,
  pdfstandard=ua-2,
  lang=en-US,
  tagging=on,
  tagging-setup={math/setup=mathml-SE,tabsorder=structure}
}
\documentclass[11pt,letterpaper]{article}
\usepackage[margin=0.68in,includefoot]{geometry}
\usepackage{newcomputermodern}
\usepackage{amsmath,amscd,mathtools}
\usepackage{tikz,adjustbox}
\providecommand{\square}{\mdlgwhtsquare}
\usepackage[hidelinks]{hyperref}
\hypersetup{
  pdftitle={Numerical Analysis PhD Exam, May 2025},
  pdfauthor={University of Florida Department of Mathematics},
  pdfsubject={Graduate mathematics examination},
  pdfdisplaydoctitle=true,
  bookmarksopen=true
}
\setcounter{secnumdepth}{0}
\setlength{\parindent}{0pt}
\setlength{\parskip}{0.28em}
\setlength{\emergencystretch}{3em}
\setlength{\leftmargini}{2.15em}
\setlength{\leftmarginii}{2.65em}
\setlength{\leftmarginiii}{3em}
\renewcommand{\labelenumi}{\textbf{\theenumi.}}
\pagestyle{empty}
\raggedbottom
\allowdisplaybreaks
\newenvironment{examproblems}[1]{%
  \begin{enumerate}%
  \setcounter{enumi}{#1}%
  \setlength{\topsep}{0.35em}%
  \setlength{\partopsep}{0pt}%
  \setlength{\itemsep}{0.48em}%
  \setlength{\parsep}{0.18em}%
}{\end{enumerate}}
\newenvironment{examsubparts}{%
  \begin{enumerate}%
  \setlength{\topsep}{0.18em}%
  \setlength{\partopsep}{0pt}%
  \setlength{\itemsep}{0.16em}%
  \setlength{\parsep}{0.1em}%
}{\end{enumerate}}
\newenvironment{examinstructions}{%
  \par\begingroup\itshape
}{%
  \par\endgroup\vspace{0.18em}
}
\makeatletter
\renewcommand\section{\@startsection{section}{1}{0pt}%
  {-1.25ex plus -.25ex minus -.1ex}%
  {0.55ex plus .1ex}%
  {\normalfont\large\bfseries}}
\renewcommand{\@maketitle}{%
  \newpage\null\vskip -1em%
  {\footnotesize\bfseries
    \href{https://www.ufl.edu/}{University of Florida}, \href{https://math.ufl.edu/}{Department of Mathematics}\par}%
  \vskip -0.5em%
  \tagstructbegin{tag=Title}%
    {\LARGE\bfseries\href{https://gma.math.ufl.edu/exams/numerical-analysis-phd/}{Numerical Analysis PhD Exam}, May 2025\par}%
  \tagstructend%
  \vskip 0.25em%
  \tagmcbegin{artifact}%
    \hrule height 1pt%
  \tagmcend%
  \vskip 1em%
}
\makeatother
\title{Numerical Analysis PhD Exam, May 2025}
\author{}
\date{}
\begin{document}
\maketitle
\thispagestyle{empty}
\begin{examinstructions}
Do 4 (four) of the first 5 (1-5) and 4 (four) of the last 5 problems (6-10).
\end{examinstructions}
\begin{examproblems}{0}
\item
Assume \(A \in \mathbb{C}^{m\times m}\).
\begin{examsubparts}
\item[\textnormal{(a)}]
Show that~\(A\) has a Schur decomposition.
\item[\textnormal{(b)}]
If~\(A\) is normal, show that~\(A\) is diagonalizable.
\end{examsubparts}
\item
Suppose~\(A\) is Hermitian positive definite.
\begin{examsubparts}
\item[\textnormal{(a)}]
Prove that each principal submatrix of~\(A\) is Hermitian positive definite.
\item[\textnormal{(b)}]
Prove that an element of~\(A\) with largest magnitude lies on the diagonal.
\item[\textnormal{(c)}]
Prove that~\(A\) has a Cholesky decomposition.
\end{examsubparts}
\item
This problem has two parts.
\begin{examsubparts}
\item[\textnormal{(a)}]
Show that \(\lVert x\rVert_\infty\) is equivalent to \(\lVert x\rVert_2\) for all \(x\in\mathbb{R}^n\). That is to find~\(C\) and~\(c\) such that \(c\lVert x\rVert_\infty\leq \lVert x\rVert_2\leq C\lVert x\rVert_\infty\), for all \(x\in\mathbb{R}^n\). Note that the constants should be determined so that the equalities hold for some nonzero \(x\in\mathbb{R}^n\).
\item[\textnormal{(b)}]
Show that \(\lVert QA\rVert_2=\lVert A\rVert_2\) if~\(Q\) is a unitary matrix.
\end{examsubparts}
\item
Assume that \(A\in\mathbb{C}^{n\times n}\) and there exists \(p\geq 1\) such that \(\lVert A\rVert_p<1\), where \(\lVert\cdot\rVert_p\) is a vector-induced matrix norm.
\begin{examsubparts}
\item[\textnormal{(a)}]
Prove that \(I-A\) is invertible.
\item[\textnormal{(b)}]
Prove that 
\[
(I-A)^{-1}=\sum_{k=0}^{\infty}A^k.
\]
\item[\textnormal{(c)}]
Prove that \(\lVert A\rVert_q\lVert A^{-1}\rVert_q\geq 1\), \(\forall 1\leq q\leq\infty\).
\item[\textnormal{(d)}]
Prove that 
\[
\frac{1}{1+\lVert A\rVert_p}\leq \lVert(I-A)^{-1}\rVert_p\leq \frac{1}{1-\lVert A\rVert_p}.
\]
\end{examsubparts}
\item
Let \(A=U\Sigma V^*\) be the singular value decomposition of \(A\in\mathbb{C}^{m\times n}\). Let \(u_j\) denote column~\(j\) of~\(U\).
\begin{examsubparts}
\item[\textnormal{(a)}]
Suppose \(\operatorname{rank}(A)=p<n<m\). Show \(\{u_1,u_2,\ldots,u_p\}\) is a basis for \(\operatorname{Col}(A)\) and \(\{u_{p+1},u_{p+2},\ldots,u_m\}\) is a basis for \(\operatorname{Null}(A^*)\).
\item[\textnormal{(b)}]
Suppose~\(A\) is full rank and \(x\neq 0\). Let \(\sigma_i\), \(i=1,\ldots,n\), be the singular values of~\(A\). Show 
\[
\sigma_1\geq \frac{\lVert Ax\rVert_2}{\lVert x\rVert_2}\geq \sigma_n>0.
\]
 If you want to use the property that \(\lVert A\rVert_2=\sigma_1\), then you must prove that it holds.
\end{examsubparts}
\item
Consider the function \(g(x)=e^{-x}\).
\begin{examsubparts}
\item[\textnormal{(a)}]
Prove that~\(g\) is a contraction on \(G=[\ln 1.1,\ln 3]\).
\item[\textnormal{(b)}]
Prove that~\(g\) maps \(G=[\ln 1.1,\ln 3]\) into \(G=[\ln 1.1,\ln 3]\).
\item[\textnormal{(c)}]
Prove that \(x_{k+1}=g(x_k)\) converges to an unique fixed point \(z\in G=[\ln 1.1,\ln 3]\) for any initial value \(x_0\in G=[\ln 1.1,\ln 3]\).
\end{examsubparts}
\item
Based on \(u_1(x)=1\), \(u_2(x)=x\), \(u_3(x)=x^2\), use Gram–Schmidt orthogonalization process to compute the three polynomials \(w_1(x)\), \(w_2(x)\), \(w_3(x)\) which are orthonormal on the interval \([0,1]\) with respect to the inner product \((f,g)=\int_0^1 f(x)g(x)\,dx\). Using these polynomials, find the best approximation in \(P^2[0,1]\) for \(f(x)=x^{1/2}\).
\item
Consider the finite difference formula 
\[
f'(t_j)=\frac{1}{12h}\bigl[f(t_j-2h)-8f(t_j-h)+8f(t_j+h)-f(t_j+2h)\bigr]+O(h^4).
\]
\begin{examsubparts}
\item[\textnormal{(a)}]
Derive this formula by using Taylor’s theorem.
\item[\textnormal{(b)}]
Derive this formula by using Lagrange polynomial representation.
\end{examsubparts}
\item
Assume the numerical quadrature for \(\hat f(\hat x)\) on \([0,1]\) is 
\[
\hat J(\hat f)=\int_0^1\hat f(\hat x)\,d\hat x\approx \hat Q(\hat f)=\sum_{j=0}^m\hat\alpha_j\hat f(\hat x_j).
\]
 Derive the numerical quadrature of \(J(f)=\int_a^b f(x)\,dx\).
\item
Consider numerically solving the initial value problem 
\[
y'(t)=f(t,y),\qquad 0<t\leq t_f,\qquad \text{with }y(0)=\eta.
\]
 Assume~\(f\) is sufficiently differentiable and let~\(h\) denote the step size. Show that all convergent members of the family of methods 
\[
y_{n+2}+(\theta-2)y_{n+1}+(1-\theta)y_n=\frac14h\bigl[(6+\theta)f_{n+2}+3(\theta-2)f_n\bigr],
\]
 parameterized by~\(\theta\), are also \(A_0\)-stable.
\end{examproblems}
\end{document}
